%% kapa: exam=1 round=31 year=2020 period=1 form=A subject=real_estate subject_ko=부동산학원론 seq=9 number=89
%% source: Q-Net 감정평가사 기출문제 (cst00309) artlSeq=5207219 file=extracted/1008/5207219/A형/A형/2020년도 제31회 감정평가사 1차 시험_1교시_A형.pdf
%% conversion: kordoc 4.24.51 -> tools/gen1.py (auto); review=proofread-sonnet; flags=none
\begin{question}{89}
다음 자료를 활용하여 거래사례비교법으로 평가한 대상토지의 감정평가액은? (단, 주어진 조건에 한함)
\begin{framedtext}
○ 대상토지: A시 B대로 30, 토지면적 200m$^2$, 제3종 일반주거지역, 주거용 토지 \par
○ 기준시점: 2020. 3. 1. \par
○ 거래사례의 내역(거래시점: 2019. 9. 1.) \par
{\small\centering
\begin{tabular}{|c|c|c|c|c|}\hline
소재지 & 용도지역 & 토지면적 & 이용상황 & 거래사례가격 \\ \hline
A시 B대로 29 & 제3종 일반주거지역 & 250m$^2$ & 주거용 & 6억원 \\ \hline
\end{tabular}\par}
○ 지가변동률(2019. 9. 1. ∼ 2020. 3. 1.): A시 주거지역은 3 \% 상승함. \par
○ 지역요인: 대상토지는 거래사례의 인근지역에 위치함. \par
○ 개별요인: 대상토지는 거래사례에 비해 8 \% 우세함. \par
○ 그 밖의 다른 조건은 동일함. \par
○ 상승식으로 계산할 것.
\end{framedtext}
\begin{choices}
  \item 531,952,000원
  \item 532,952,000원
  \item 533,952,000원
  \item 534,952,000원
  \item 535,952,000원
\end{choices}
\end{question}
\answer{3}
